\documentclass[11pt]{article}
\usepackage[a4paper,margin=25mm]{geometry}
\usepackage{amsmath,amssymb,amsthm}
\usepackage[hidelinks]{hyperref}
\emergencystretch=2em
\newtheorem{theorem}{Theorem}
\newcommand{\R}{\mathbb R}
\title{Identical Asymptotics with Different Finite Iterates:\\
A Proof of Conjecture 00000006090}
\author{GitHub: gaochengzhecpu}
\date{}
\begin{document}
\maketitle
\begin{abstract}
For any prescribed finite iteration, two gradient-descent step schedules
on the same quadratic objective, with the same initial value, can produce
different errors and then rejoin exactly one step later. Both sequences
converge to the same minimizer and have asymptotic successive-error ratio
$1/2$. In fact, their entire tails coincide. The objective derivative,
actual iterative algorithms, explicit trajectories, exact error gap,
convergence, and quantified existence claim are verified in Lean.
\end{abstract}

\paragraph{Reading of the source.}
Both versions describe the desired distinction as different
``finite-scale behavior'' or ``pre-asymptotics,'' despite the same limit
and asymptotic convergence rate. We realize this using actual errors
of iterative algorithms at finite iteration indices. The source does
not define a separate operator or numerical invariant called the
finite-scale spectrum, and none is introduced here. We make the
convergence rate precise by the standard successive-error ratio and
also prove the stronger property that the trajectories eventually
coincide exactly. Thus the distinction cannot depend on choosing a
weak interpretation of equal asymptotic rates.

\begin{theorem}\label{thm:main}
Let $f(x)=x^2/2$ on $\R$ and start gradient descent at $x_0=y_0=1$.
For every integer $m\geq0$, there exist two step-size schedules
$a_n,b_n\in(0,1)$ such that
\[
 x_{n+1}=x_n-a_nf'(x_n),\qquad
 y_{n+1}=y_n-b_nf'(y_n)
\]
and the following hold:
\begin{enumerate}
\item $x_n=y_n$ for $0\leq n\leq m$;
\item at the prescribed next iteration,
\[
 |y_{m+1}|-|x_{m+1}|=2^{-(m+2)}>0;
\]
\item $x_n=y_n$ for every $n\geq m+2$;
\item both sequences converge to the unique minimizer zero and
\[
 \lim_{n\to\infty}\frac{|x_{n+1}|}{|x_n|}
 =\lim_{n\to\infty}\frac{|y_{n+1}|}{|y_n|}=\frac12.
\]
\end{enumerate}
Consequently, the algorithms have the same asymptotic convergence
rate and limit, but different finite-scale behavior.
\end{theorem}

\begin{proof}
The objective is nonnegative, vanishes exactly at zero, and has
derivative $f'(x)=x$. Define
\[
 a_n=\frac12,\qquad
 b_n=\begin{cases}
 \frac14,&n=m,\\[2pt]
 \frac23,&n=m+1,\\[2pt]
 \frac12,&\text{otherwise}.
 \end{cases}
\]
Each step lies strictly between zero and one. Gradient descent becomes
$x_{n+1}=(1-a_n)x_n$ and $y_{n+1}=(1-b_n)y_n$. Thus
\begin{equation}\label{eq:trajectories}
 x_n=2^{-n},\qquad
 y_n=\begin{cases}
 \frac32\,2^{-n},&n=m+1,\\[2pt]
 2^{-n},&n\neq m+1.
 \end{cases}
\end{equation}
To check the second formula directly, the sequences follow the same
updates up to index $m$, where both equal $2^{-m}$. The changed step
$1/4$ then gives
\[
 y_{m+1}=\frac34\,2^{-m}
 =\frac32\,2^{-(m+1)}.
\]
The next changed step compensates exactly:
\[
 y_{m+2}=\left(1-\frac23\right)y_{m+1}
 =\frac13\cdot\frac34\,2^{-m}=2^{-(m+2)}=x_{m+2}.
\]
All later steps again equal $1/2$, so the trajectories remain identical.
These calculations prove \eqref{eq:trajectories}; alternatively its
substitution in the recurrence proves it by induction.

All displayed iterates are positive. Therefore at $m+1$ the exact
error difference is
\[
 |y_{m+1}|-|x_{m+1}|
 =\left(\frac32-1\right)2^{-(m+1)}=2^{-(m+2)}.
\]
The geometric sequence $2^{-n}$ tends to zero. By eventual equality,
so does $y_n$. The successive-error ratio of $x_n$ equals $1/2$
at every index. For $n\geq m+2$, both $y_n$ and $y_{n+1}$ equal
the corresponding $x$ iterates, so its ratio is also exactly $1/2$.
This proves all assertions. Since $m$ is arbitrary, the differing
finite iteration can be placed as late as desired, while the tails
and therefore the asymptotic behavior still coincide exactly.
\end{proof}

\section*{Formal verification and scope}
The file \texttt{lean/Main.lean} uses Lean 4.19.0 with Mathlib commit
\begin{center}\texttt{c44e0c8ee63ca166450922a373c7409c5d26b00b}.\end{center}
The objective is a real function and its derivative is verified using
Mathlib's \texttt{HasDerivAt}. The two algorithms are actual recursively
defined gradient-descent iterates using Mathlib's \texttt{deriv}, not
numerical sequences named as algorithms without proving their update
rule. The same objective and initial value are used for every $m$.

Lean proves the explicit trajectories and the step bounds $(0,1)$.
It proves equality of the prefixes and tails, the exact error gap,
and the finite error discrepancy. Convergence to zero is a filter
limit. The predicate \texttt{HasRate} gives the limit of the ratio of
successive absolute errors. It is proved for both sequences.
Their eventual equality is also formalized directly.

The final theorem is
\begin{center}\texttt{Conjecture6090.conjecture\_6090\_true}.\end{center}
It proves \texttt{SeparationClaim}: there exists a single differentiable
objective with unique zero minimum, such that for every $m$ there exist
two positive step schedules and their gradient-descent trajectories
with the stated limit, rate, finite error discrepancy, common prefix,
and identical tail. This is stronger than the source's existential
claim. All mathematical ingredients needed for the theorem are
formalized; no unspecified spectral object is substituted for the
actual finite error behavior.

The audited main, derivative, trajectory, gap, and rate theorems use
only \texttt{propext}, \texttt{Classical.choice}, and \texttt{Quot.sound}.
This submission was developed by a delegated AI agent with a
self-review; the coordinating agent reviews it separately. No
external independent review is claimed.

\section*{Source and reproduction}
The exact bilingual source is included as \texttt{SOURCE.md}, from
upstream commit
\begin{center}\texttt{b355e42906e6eaf52bc7e9d005e0ef02d8aa7400}.\end{center}
The source hash and retrieval record are in
\begin{center}\texttt{verification/SOURCE\_PROVENANCE.md}.\end{center}
The proof is elementary and self-contained. With the pinned
dependencies available, run in the \texttt{lean} directory:
\begin{center}\texttt{lake env lean -DwarningAsError=true Main.lean}\end{center}
and \texttt{lake build} for a fresh build. Compile \texttt{main.tex}
with Tectonic or a compatible LaTeX engine.
\end{document}
